\section{总结与延伸}

本讲可以被压缩为一条“曲线解释链”。第一，language-model loss 随 compute、data 与 parameters 平滑改善；第二，task metric 由于 chance baseline、exact match 和多子技能瓶颈，可能长时间平坦后上升；第三，这个阈值会被 model family、data quality、architecture、fine-tuning 与 measurement axis 移动；第四，CoT 改变 elicitation interface，让较大模型用中间 token 展开组合计算；第五，self-consistency 再用多路径 inference budget 聚合答案；第六，所有结论都必须与成本、benchmark construction 和历史设置一起报告。

\subsection{一张术语消化表}

下面的术语经常在一句话里被混用。把它们拆开，才能判断某篇论文究竟增加了训练资源、模型能力、提示信息还是推理时搜索。

\begin{longtable}{L{0.19\textwidth}L{0.32\textwidth}L{0.40\textwidth}}
\toprule
\textbf{术语} & \textbf{解决的问题} & \textbf{核心机制与课程关系} \\
\midrule
\endhead
Scaling law & 预测平均 loss 怎样随资源变化 & 在受控模型族上拟合 compute、data、parameters 与 loss 的平滑关系，是 emergence 的对照基线。 \\
Emergent ability & 描述小模型 metric 不可见、大模型可见的任务能力 & 依赖 chance/human baseline、prompt、metric 与 model family，不等于已证明离散内部模块。 \\
Few-shot prompting & 不更新权重地告诉模型任务格式 & 把 exemplar 放入 context，让 next-token model 归纳输入输出映射。 \\
Instruction fine-tuning & 让模型更稳定遵循自然语言任务 & 更新 weights，可将已知 desired behavior 提前诱导到更小模型。 \\
RLHF & 让输出更符合人类 preference & 使用 comparison/reward signal 更新模型；与仅改变 context 的 CoT 不同。 \\
Chain-of-thought & 为多步任务提供可生成的中间状态 & 在示例中写 reasoning trace，让模型按 token 顺序展开 decomposition。 \\
Perplexity & 衡量 held-out next-token prediction & $\exp(\text{cross-entropy})$；可作有效模型质量 proxy，但不是下游能力总分。 \\
Self-consistency & 降低单条 CoT 偶然走错的风险 & 采样多条 reasoning paths，对 final answer 投票，以 inference cost 换 accuracy。 \\
Temperature & 控制 next-token sampling 的随机性 & 改变 logit softmax 的平滑程度，为 self-consistency 提供 path diversity。 \\
\bottomrule
\end{longtable}

\subsection{如何用本讲框架评审新的 Scaling 声明}

面对新的 ``emergent'' 图，先要求原作者公开每个 checkpoint、prompt、baseline、seed 与 confidence interval，再检查 metric 是否把连续质量变成硬阈值；面对新的 reasoning method，要求同时报告 base model、generated tokens、sample count、latency 与 external tools；面对跨模型家族结论，则必须询问 data 和 post-training 是否也改变。这样才能把“看起来突然”转化为可重复的工程事实。

\begin{warningbox}{三种最常见的过度结论}
第一，从几个稀疏 checkpoint 推断精确 phase transition；第二，把参数量与数据、架构、post-training 的共同变化全归因于参数；第三，把可读 CoT 当作模型内部真实思维。三者都比课堂证据更强，应在写作和评审中主动降级。
\end{warningbox}

\subsection{拓展阅读}

以下资料均为课堂日期前已经公开、且直接构成本讲证据链的一手来源。阅读时建议先复现一张曲线的 axes、baseline 与 prompt，再比较论文作者对机制的表述强度。

\begin{itemize}
  \item \href{https://arxiv.org/abs/2001.08361}{Kaplan et al., Scaling Laws for Neural Language Models}：平滑 loss scaling 的经典对照。
  \item \href{https://arxiv.org/abs/2206.07682}{Wei et al., Emergent Abilities of Large Language Models}：任务与 prompting-technique 的涌现目录。
  \item \href{https://arxiv.org/abs/2211.02011}{Wei, Tay, and Le, Inverse Scaling Can Become U-Shaped}：局部负趋势为何可能反转。
  \item \href{https://arxiv.org/abs/2201.11903}{Wei et al., Chain-of-Thought Prompting Elicits Reasoning in Large Language Models}：CoT 主论文。
  \item \href{https://arxiv.org/abs/2210.09261}{Suzgun et al., Challenging BIG-Bench Tasks and Whether Chain-of-Thought Can Solve Them}：BBH、任务异质性与 scale threshold。
  \item \href{https://arxiv.org/abs/2210.03057}{Shi et al., Language Models are Multilingual Chain-of-Thought Reasoners}：MGSM 与跨语言组合性。
  \item \href{https://arxiv.org/abs/2203.11171}{Wang et al., Self-Consistency Improves Chain of Thought Reasoning in Language Models}：多路径采样与多数投票。
  \item \href{https://arxiv.org/abs/2210.11416}{Chung et al., Scaling Instruction-Finetuned Language Models}：task、model 与 CoT data scaling 的后训练证据。
\end{itemize}

最终，本讲最重要的训练不是记住某个参数阈值，而是养成一种曲线纪律：先问测量什么，再问改变什么，最后才问能否外推。Scaling 让模型进入新的可达区域，数据和后训练移动边界，CoT 与 self-consistency 改变我们探索区域的方式；真正可靠的结论，必须同时说明能力、接口、成本与不确定性。

\subsection{课后自检}

读完本讲后，读者应能不依赖“bigger is better”口号，独立拆解一项新的能力声明。下面的问题可作为论文阅读、实验设计和产品评估的最低检查表。

\begin{importantbox}{六个自检问题}
\begin{enumerate}
  \item 为什么 smooth cross-entropy 改善可以与 exact-match accuracy 的突然上升同时发生？
  \item 参数量、训练 FLOPs、数据量与 perplexity 分别混入了哪些假设？
  \item 怎样用 controlled ablation 区分 better data 与 model-family confounder？
  \item CoT 改变了 weights、context、generated-token budget 中的哪一项？
  \item Self-consistency 在什么错误相关结构下会失败，即使 sample count 很大？
  \item 一张 emergence 图还缺哪些 baseline、scale points 或 cost information 才足以支持工程决策？
\end{enumerate}
\end{importantbox}

\end{document}
